\documentclass[11pt,a4paper]{article}

\usepackage[margin=2.5cm]{geometry}
\usepackage{amsmath,amssymb,amsthm,mathtools}
\usepackage{mathrsfs}
\usepackage{booktabs}
\usepackage{microtype}
\usepackage[hidelinks]{hyperref}

\newtheorem{proposition}{Proposition}
\newtheorem{observation}{Observation}
\newtheorem{problem}{Problem}
\setcounter{problem}{15}
\newtheorem{remark}{Remark}
\newtheorem{erratumthm}{Erratum}

\newcommand{\E}{\mathbb E}
\newcommand{\D}{\mathscr D}

\title{\textbf{Erratum: ``Bulk is Gauge'' Was a Quantifier Error}\\
\large The Joint Bulk Floor is Nonzero, and Larger Than the Ghost Mode}
\author{}
\date{August 29, 2026}

\begin{document}
\maketitle

\begin{abstract}
Observation~1 of the preceding note (``Bulk is gauge'') claimed that every
tested bulk antidiagonal cell has minimal achievable spread exactly zero,
and concluded that bulk routing carries no forced history-dependence. This
conflated two different statements: that every cell \emph{individually}
admits a face point with zero spread, versus that \emph{one} face point
gives every cell zero spread simultaneously. The first is what was tested;
the second is what the conclusion needed. Running the correct joint test --
minimizing the worst-case spread across all bulk cells under a single
shared point of the near-optimal face -- gives a floor of
$1.475854\times10^{-2}$ at depth 3, an order of magnitude \emph{larger}
than the previously reported ghost-cell floor of $1.557\times10^{-3}$.
Bulk routing is not gauge. Six of the seven tested cells sit at exactly the
same spread value at the joint optimum, indicating a shared coupling
constraint rather than seven independent phenomena.
\end{abstract}

\section{The error}

Observation~1 of the prior note stated: ``At depth 4, for every tested
$(s,x)$... $\mathrm{spread}^\star(s,x)=0$... Every bulk routing pattern
previously observed to vary across histories was solver gauge, not forced
structure.'' The computation behind the first sentence is correct: each
cell, tested in isolation, does admit a face point with zero spread. The
second sentence does not follow from it. Symbolically, what was shown is
\[
(\forall (s,x))\,(\exists z_{(s,x)}\in F^\star)\ \ \mathrm{spread}(s,x;z_{(s,x)})=0,
\]
where $F^\star$ is the near-optimal face and $z_{(s,x)}$ may depend on
$(s,x)$. What ``bulk is gauge'' requires is the quantifier-reversed
statement
\[
(\exists z\in F^\star)\,(\forall (s,x))\ \ \mathrm{spread}(s,x;z)=0,
\]
a single $z$ working for every cell at once. These are not equivalent, and
only the second licenses the conclusion that was drawn.

\section{The corrected test}

\begin{proposition}[Joint face-spread LP]
For observables $q=1,\dots,m$ (here, the tested bulk cells), the quantity
\[
\mathrm{jointfloor}^\star := \min_{z\in F^\star}\ \max_{q}\Bigl(\max_h L_{q,h}(z) - \min_h L_{q,h}(z)\Bigr)
\]
is computable by one linear program: introduce $(U_q,L_q)$ per observable
and a shared $t$ with $U_q-L_q\le t$ for every $q$, and minimize $t$ subject
to the same face constraint $|c^\top z-v^\star|\le\tau$ used throughout.
$\mathrm{jointfloor}^\star=0$ correctly certifies that bulk is gauge;
$\mathrm{jointfloor}^\star>0$ certifies that it is not, regardless of what
the per-cell minima were.
\end{proposition}

\section{Corrected result}

\begin{observation}[Retraction]
At depth 3, $M=8$, testing the seven cells $s\in\{2,3\}$ (all admissible
$x$), every individual cell has $\mathrm{spread}^\star(s,x)=0$, reproducing
the earlier (correct, but insufficient) per-cell computation. The joint
floor is
\[
\mathrm{jointfloor}^\star = 1.475854\times10^{-2},
\]
five orders of magnitude above the numerical tolerance. \textbf{Observation
1 of the prior note is retracted}: bulk routing is not gauge. No single
point of the near-optimal face removes history-dependence from every bulk
cell at once.
\end{observation}

\begin{observation}[A shared bottleneck, not seven independent facts]
At the joint optimizer, six of the seven cells --
$(s,x)\in\{(2,0),(2,1),(2,2),(3,1),(3,2),(3,3)\}$ -- have spread exactly
equal to the joint floor, $1.475854\times10^{-2}$; the seventh, $(3,0)$, has
a strictly smaller spread, $4.983532\times10^{-3}$, at that same point.
Six unrelated cells landing on the identical numerical value at the
optimizer is not expected of seven independent obstructions; it is the
signature of one shared coupling constraint (most plausibly a conservation
or flow-consistency relation linking these cells' routing) that the
optimizer cannot relax without paying elsewhere. Identifying that
constraint explicitly is more promising than treating the seven cells as
separate data points.
\end{observation}

\begin{remark}[Larger than the mode we were chasing]
The corrected bulk floor, $1.48\times10^{-2}$, is roughly ten times the
ghost-cell floor, $1.557\times10^{-3}$, reported as the main finding of the
prior note. The ghost cell was not wrong -- it remains independently
confirmed at two depths and two tolerances -- but it is evidently not the
dominant source of forced structure. The bulk coupling is.
\end{remark}

\section{What remains to be done}

The depth-4 joint bulk floor did not complete within the time budget
available at the time of writing (the corresponding single-cell tests at
depth 4 already ran to several hundred seconds each; the joint LP couples
all cells into one larger program). Obtaining it, and comparing to the
depth-3 value here, is the immediate next step: if the joint floor is
roughly stable across depth (as the ghost cell's floor was, Section~4 of
the prior note), that is further evidence for a marginal, scale-stable
mode; if it grows with depth, that is evidence for a relevant, dangerous
one, in the renormalization language used elsewhere in this project.

\begin{problem}[Identify and track the joint coupling]
Determine the linear relation among $z$ responsible for pinning six of the
seven tested cells to a common spread value, and track
$\mathrm{jointfloor}^\star$ across depth (at minimum depth 4, ideally depth
5) to determine whether this mode is bounded, shrinking, or growing with
scale.
\end{problem}

\section{Scope}

This note corrects a specific, identified overreach in a previously
circulated result and reports the outcome of the correct test in its
place. It does not yet supply the depth-4 joint floor, does not identify
the specific coupling constraint responsible for the six-way tie, and does
not change the status of the ghost-cell finding, which stands independently
of this correction.

\end{document}
